\originalsection{18.315：作业4（原卷 Fall 2005）}
本组收录原题1、4、6, 共四个任务单元.
原题2只给 Stanley《计数组合学》卷I题2.11,
原题3只给 Bollobás《现代图论》第376页题45、46;
官方题卷没有这些外书题面, 故只登记缺项, 不按编号猜题或再造题面.

\begin{exercise}\label{mit:18315-hw4-q1}
\textbf{18.315, 原卷 Fall 2005, 作业4, 题1.}
用矩阵树定理求完全 $r$ 部图 $K_{n_1,\ldots,n_r}$ 的生成树数,
其中各部分大小为正整数 $n_1,\ldots,n_r$.
\end{exercise}
\sourceof{官方 HW4 第1页题1; 编者独立解答}
\begin{solution}
记 $N=\sum_i n_i$. 当 $r\ge2$ 时, 图连通, 答案为
\[
\tau(K_{n_1,\ldots,n_r})=N^{r-2}\prod_{i=1}^r(N-n_i)^{n_i-1}.
\]
求其 Laplace 矩阵 $L$ 的非零特征值.
在第 $i$ 部分内取坐标和为0、其他部分为0的向量,
同部没有边, 得到特征值 $N-n_i$, 重数 $n_i-1$.
剩余子空间由各部分上为常数 $c_i$ 的向量组成,
$L$ 在该部分的坐标为 $Nc_i-\sum_jn_jc_j$.
全常数向量给特征值0, 而 $\sum_jn_jc_j=0$ 的 $r-1$ 维子空间给特征值 $N$.

若连通图的非零 Laplace 特征值为 $\lambda_2,\ldots,\lambda_N$,
则 $\det(L+tI)=t\prod_{j=2}^N(t+\lambda_j)$ 的 $t$ 系数等于它们的乘积.
按行列式展开, 该系数又是全部 $N$ 个主余子式之和,
矩阵树定理说每个主余子式都等于 $\tau$.
所以 $N\tau=\prod_{j=2}^N\lambda_j$, 代入上面的谱即得到所示公式.
退化情形 $r=1$: $N=1$ 时空生成树计为1, $N>1$ 时图无边且不连通, 树数为0.
\end{solution}

\begin{exercise}\label{mit:18315-hw4-q4}
\textbf{18.315, 原卷 Fall 2005, 作业4, 题4.}
记 $C(k,n,q)$ 为 $k\times n$ 网格图的正常 $q$ 染色数.
证明: (a) 固定 $k,n$ 时, 可用关于 $\log q$ 的多项式时间计算它;
(b) 固定 $k,q$ 时, 可用关于 $n$ 的多项式时间计算它.
\end{exercise}
\sourceof{官方 HW4 第1页题4; 编者独立解答, 时间指整数位运算时间}
\begin{solution}
(a) 固定的图有 $v=kn$ 个顶点. 删除--收缩恒等式可离线计算它的染色多项式
$P_G(z)=\sum_{j=0}^v a_jz^j$, 系数只依赖这一个固定图.
输入 $q$ 后用 Horner 法逐次计算, 运算次数为常数,
中间整数的位数为 $O(\log(q+2))$, 因为多项式次数与系数固定.
普通整数加法和乘法的位复杂度均为多项式, 因而总时间亦为多项式.
写 $\log(q+2)$ 同时涵盖 $q=0,1$ 的输入, 与原题通常 $q\ge2$ 的表述等价.

(b) 把一列的 $k$ 个颜色作为状态. 合法状态集
$\mathcal S=\{(c_1,\ldots,c_k)\in[q]^k:c_i\ne c_{i+1}\}$ 固定不变.
令转移矩阵 $M$ 的 $(s,t)$ 元素为1当且仅当对每个位置 $i$ 有 $s_i\ne t_i$,
否则为0. 这恰保证相邻两列的水平边也正常.
于是
\[
C(k,n,q)=\boldsymbol1^{\mathsf T}M^{n-1}\boldsymbol1.
\]
由第一列各状态计数1开始, 逐列乘矩阵即可.
状态数至多 $q^k$, 是固定常数; 各步计数不超过 $q^{kn}$,
位数为 $O(n)$（$q\ge2$）, 因此 $n-1$ 次固定大小整数矩阵乘法的位时间为多项式.
$q=0,1$ 的情形可直接检查, 或由空/单状态矩阵处理.
原题要求关于 $n$ 的多项式时间, 没有要求关于其二进制长度 $\log n$ 的时间;
一般答案本身就有 $\Theta(n)$ 位.
\end{solution}

\begin{exercise}\label{mit:18315-hw4-q6}
\textbf{18.315, 原卷 Fall 2005, 作业4, 题6.}
设 $A(m)$ 是恰有 $m$ 条边的图所能拥有的最大生成树数.
求非平凡上下界, 改进原题所列的 $2^m$ 上界及直接取一个完整大完全图的下界.
\end{exercise}
\sourceof{官方 HW4 第1页题6; 编者独立解答, 所得界不宣称最优}
\begin{solution}
若图有 $v$ 个顶点, 每棵生成树恰选 $v-1$ 条不同边,
故其数不超过 $\binom m{v-1}$, 不连通图的树数则为0.
在二项式系数中间项最大, 因而
\[
A(m)\le\binom m{\lfloor m/2\rfloor}
\le\frac{2^m}{\sqrt{\lfloor m/2\rfloor+1}}.
\]
后一界可直接证出, 无需使用渐近公式.
当 $m=2t$ 时,
\[
\frac{\binom{2t}t}{4^t}
=\prod_{j=1}^t\left(1-\frac1{2j}\right)
\le\exp\!\left(-\frac12\sum_{j=1}^t\frac1j\right)
\le(t+1)^{-1/2};
\]
最后用 $\sum_{j=1}^t1/j\ge\log(t+1)$.
奇数 $m=2t+1$ 时,
$\binom{2t+1}t=\frac{2t+1}{t+1}\binom{2t}t\le2\binom{2t}t$,
得到同一形式的界.

下界取 $a=\lfloor m/10\rfloor$ 份 $K_5$, 只把它们的一个指定顶点合并为同一割点,
再接一条有 $m-10a$ 条边的路径. 全图简单、连通且恰有 $m$ 条边.
每棵生成树在每个 $K_5$ 块上的限制必须是该块的生成树,
反之各块生成树的并加上路径各边恰为全图生成树;
所以由 Cayley 公式,
\[
A(m)\ge125^{\lfloor m/10\rfloor}.
\]
这给出指数下界, 而随 $m$ 增大的单一 $K_v$ 只有
$\log\tau=(v-2)\log v=O(\sqrt m\log m)$.
本解答分别从固定大小块的乘法结构和生成树的固定边数改进了原题两界,
没有声称计算出了最优的 $A(m)$.
\end{solution}
